\begin{figure}[htbp]
  \centering
  \begin{tikzpicture}[x=7mm,y=-7mm,alt={A three by five source contains rows zero through four, five through nine, and ten through fourteen. Its five by three transpose contains rows zero five ten; one six eleven; two seven twelve; three eight thirteen; four nine fourteen. Shading covers source columns zero through three in all three rows and the corresponding first four complete destination rows. This is the first clipped four by four source tile, with twelve cells.}]
    \foreach \i in {0,1,2} {\foreach \j in {0,1,2,3} {\fill[black!12] (\j,\i) rectangle ++(1,1);}}
    \foreach \i in {0,1,2,3} {\fill[black!12] (7,\i) rectangle ++(3,1);}
    \foreach \x in {0,...,5} {\draw (\x,0)--(\x,3);}
    \foreach \y in {0,...,3} {\draw (0,\y)--(5,\y);}
    \foreach \x in {7,...,10} {\draw (\x,0)--(\x,5);}
    \foreach \y in {0,...,5} {\draw (7,\y)--(10,\y);}
    \foreach \i in {0,1,2} {\foreach \j in {0,...,4} {\pgfmathtruncatemacro{\n}{5*\i+\j}\node[font=\small] at (\j+.5,\i+.5) {\n};}}
    \foreach \i in {0,...,4} {\foreach \j in {0,1,2} {\pgfmathtruncatemacro{\n}{5*\j+\i}\node[font=\small] at (7+\j+.5,\i+.5) {\n};}}
    \node[anchor=south,font=\small\bfseries] at (2.5,-.2) {Source: $3\times5$};
    \node[anchor=south,font=\small\bfseries] at (8.5,-.2) {Output: $5\times3$};
    \draw[->,thick] (5.35,1.4)--(6.65,1.4);
  \end{tikzpicture}
  \caption{Transpose changes row width as well as coordinates. The source values name their original linear indices. Shading identifies the first clipped tile; the remaining source column becomes the final output row.}
  \label{fig:revision-locality-mapping}
\end{figure}
